%═══════════════════════════════════════════════════════════
% APPENDIX A: NOTATION TABLE
% 
% Prerequisites: All chapters (references notation used throughout)
% Provides: Comprehensive symbol reference
% Labels created: app:notation
% Page budget: 1-2 pages (longtable handles page breaks)
% Status: COMPLETE through Chapter 3, plus S4.1. Ch3 contributes wc(T),
%   the upper-bound vector u, and the three S3.4 symbols pi, Pi and
%   [z]^+. S4.1 contributes wcr(T). Proof-local symbols are defined at
%   the point of use and stay out of this table by design: v_{i,j},
%   a_j, b_j, r_j, T_P, W, Q, w(P) in §3.3.1; a_i, b_i, v_{i,j},
%   L_{i,j}, p in §3.3.2; f_j, g, c^*, c^pi and phi_T in §3.4.
%
% REQUIRES: \usepackage{longtable} in main.tex preamble
%           Add after \usepackage{booktabs}
%═══════════════════════════════════════════════════════════

\chapter{Notation}\label{app:notation}

This appendix provides a reference for the mathematical notation used throughout the thesis.
Symbols are grouped by category with the section of first use indicated.

\begin{longtable}{@{}l>{\raggedright\arraybackslash}p{0.56\textwidth}l@{}}
\toprule
\textbf{Symbol} & \textbf{Meaning} & \textbf{First Use} \\
\midrule
\endfirsthead

\toprule
\textbf{Symbol} & \textbf{Meaning} & \textbf{First Use} \\
\midrule
\endhead

\midrule
\multicolumn{3}{r}{\textit{Continued on next page}} \\
\endfoot

\bottomrule
\endlastfoot

\multicolumn{3}{l}{\textit{Graph Notation}} \\
$G = (V, E)$ & Undirected graph with vertex set $V$, edge set $E$ & \S\ref{sec:graph-notation} \\
$n = |V|$, $m = |E|$ & Number of vertices and edges & \S\ref{sec:graph-notation} \\
$\cut{X}$ & Cut induced by vertex set $X \subseteq V$ & \S\ref{sec:graph-notation} \\
\midrule

\multicolumn{3}{l}{\textit{Trees and Costs}} \\
$\cT$ & Set of all spanning trees of $G$ & \S\ref{sec:graph-notation} \\
$T$ & Spanning tree, $T \in \cT$ & \S\ref{sec:graph-notation} \\
$c \colon E \to \R$ & Edge cost function & \S\ref{sec:graph-notation} \\
$c(T)$ & Cost of tree $T$, i.e., $\sum_{e \in T} c_e$ & \S\ref{sec:graph-notation} \\
$\MSTcost{c}$ & Minimum spanning tree cost under $c$ & \S\ref{sec:graph-notation} \\
$T^*$ & A minimum spanning tree, $c(T^*) = \MSTcost{c}$ & \S\ref{sec:graph-notation} \\
$C_f$ & Fundamental cycle of non-tree edge $f$ & \S\ref{sec:graph-notation} \\
$\cut{X_e}$ & Fundamental cut of tree edge $e$ & \S\ref{sec:graph-notation} \\
\midrule

\multicolumn{3}{l}{\textit{Uncertainty Models}} \\
$\Scenarios$ & Uncertainty set (discrete, interval, or budgeted) & \S\ref{sec:uncertainty} \\
$K$ & Number of discrete scenarios & \S\ref{sec:uncertainty} \\
$\cs{k}$ & Scenario $k$ cost vector, $k \in \{1, \ldots, K\}$ & \S\ref{sec:uncertainty} \\
$[\ell_e, u_e]$ & Interval cost bounds for edge $e$ & \S\ref{sec:uncertainty} \\
$\ell_e, u_e$ & Lower and upper bounds on edge cost (interval model) & \S\ref{sec:uncertainty} \\
$\Scenarios^{\Gamma}$ & Budgeted uncertainty set & \S\ref{sec:uncertainty} \\
$\Gamma$ & Deviation budget (budgeted model) & \S\ref{sec:uncertainty} \\
$\bar{c}_e$ & Nominal cost of edge $e$ (budgeted model) & \S\ref{sec:uncertainty} \\
$\hat{c}_e$ & Maximum deviation of edge $e$ (budgeted model) & \S\ref{sec:uncertainty} \\
$\delta_e$ & Fraction of its maximum deviation taken by edge $e$ & \S\ref{sec:uncertainty} \\
$\R^{|E|}_{\geq 0}$ & Non-negative cost vectors, required by all three models & \S\ref{sec:uncertainty} \\
$u = (u_e)_{e \in E}$ & Upper-bound cost vector, to which the interval model reduces & \S\ref{sec:mm-extremal} \\
$c^{T}$ & Scenario charging $u_e$ on $E(T)$ and $\ell_e$ elsewhere, at which the maximum regret of $T$ is attained & \S\ref{sec:regret-extremal} \\
$c^{\mathrm{av}}$ & Midpoint cost vector, $c^{\mathrm{av}}_e = \tfrac{1}{2}(\ell_e + u_e)$ & \S\ref{sec:regret-approx-interval} \\
$\pi$ & Level above which deviations are charged (budgeted reformulation) & \S\ref{sec:mm-budgeted} \\
$\Pi$ & Candidate levels $\{0\} \cup \{\hat{c}_e : e \in E\}$ & \S\ref{sec:mm-budgeted} \\
$[z]^{+}$ & Positive part, $\max\{z, 0\}$ & \S\ref{sec:mm-budgeted} \\
\midrule

\multicolumn{3}{l}{\textit{Robust Objectives}} \\
$\Regret{T}{c}$ & Regret of $T$ under $c$: $c(T) - \MSTcost{c}$ & \S\ref{sec:uncertainty} \\
$\wc{T}$ & Worst-case cost of tree $T$: $\max_{c \in \Scenarios} c(T)$ & \S\ref{sec:mm-formulation} \\
$\wcr{T}$ & Maximum regret of tree $T$: $\max_{c \in \Scenarios} \Regret{T}{c}$ & \S\ref{sec:regret-definition} \\
\midrule

\multicolumn{3}{l}{\textit{Complexity}} \\
$\mathsf{P}$ & Polynomial-time complexity class & \S\ref{sec:complexity} \\
$\mathsf{NP}$ & Nondeterministic polynomial-time class & \S\ref{sec:complexity} \\
$\mathsf{NP}$-hard & At least as hard as $\mathsf{NP}$-complete problems & \S\ref{sec:complexity} \\
Weakly $\mathsf{NP}$-hard & $\mathsf{NP}$-hard with a pseudo-polynomial algorithm & \S\ref{sec:complexity} \\
Strongly $\mathsf{NP}$-hard & $\mathsf{NP}$-hard even with poly-bounded numeric values & \S\ref{sec:complexity} \\
Pseudo-polynomial & Algorithm polynomial in input size and numeric magnitudes & \S\ref{sec:complexity} \\
$\alpha$-approximation & Algorithm guaranteeing $\mathrm{ALG} \leq \alpha \cdot \OPT$ & \S\ref{sec:complexity} \\
$\varepsilon$ & Tolerance parameter for approximation schemes & \S\ref{sec:complexity} \\
PTAS & Polynomial-time approximation scheme & \S\ref{sec:complexity} \\
FPTAS & Fully polynomial-time approximation scheme & \S\ref{sec:complexity} \\

\end{longtable}

% END OF APPENDIX A